\documentclass[10pt]{article}
\usepackage[a4paper,margin=0.86in]{geometry}
\usepackage{amsmath,amssymb,amsthm,mathtools,enumitem,xcolor,hyperref}
\hypersetup{colorlinks=true,linkcolor=blue!55!black,citecolor=blue!55!black,urlcolor=blue!55!black}
\setlength{\emergencystretch}{3em}
\setlength{\parindent}{0pt}
\setlength{\parskip}{0.52em}
\setlist{leftmargin=1.6em,itemsep=0.22em,topsep=0.25em}
\sloppy

\newtheorem{theorem}{Theorem}[section]
\newtheorem{proposition}[theorem]{Proposition}
\newtheorem{corollary}[theorem]{Corollary}
\theoremstyle{definition}
\newtheorem{definition}[theorem]{Definition}
\newtheorem{convention}[theorem]{Convention}
\newtheorem{remark}[theorem]{Remark}

\newcommand{\Fq}{\mathbf F_q}
\newcommand{\Fbar}{\overline{\mathbf F}_q}
\newcommand{\Ql}{\mathbf Q_\ell}
\newcommand{\Qlb}{\overline{\mathbf Q}_\ell}
\newcommand{\Dbc}{D^b_c}
\newcommand{\WeilD}{D^b_{\mathrm W,c}}
\newcommand{\WtLe}{\operatorname{WeightLe}}
\newcommand{\WtGe}{\operatorname{WeightGe}}
\newcommand{\Fr}{\operatorname{Fr}}
\newcommand{\Spec}{\operatorname{Spec}}
\newcommand{\gr}{\operatorname{gr}}
\newcommand{\RHom}{R\mathcal Hom}

\begin{document}
\begin{center}
{\large\bfseries Weil Structures, Frobenius, and Weights\par}
\medskip
{\normalsize Frobenius conventions, coefficient categories, and weight calculus\par}
\end{center}

\begin{abstract}
This note defines every Frobenius and weight datum used in the cohomological
Weil conjectures.  It distinguishes arithmetic field Frobenius, geometric
field Frobenius, and scheme Frobenius; constructs the derived Weil category
by localizing strict equivariant complexes; and constructs algebraically
closed coefficients from finite coefficient fields.  Punctual purity,
mixedness, the BBD upper and lower predicates, scalar and base-field change,
and the monodromy filtration are then stated as a sequence of definitions
and theorems.  All derived categories are ordinary triangulated categories.
\end{abstract}
\tableofcontents

\section{Geometric and coefficient conventions}

Let \(k=\Fq\) have characteristic \(p\), let \(\ell\ne p\), and fix an
algebraic closure \(\bar k=\Fbar\).  Let \(X_0\) be a separated scheme of
finite type over \(k\), and put
\[
                         X=X_0\times_k\bar k.
\]
All sheaves in this note are sheaves on the pro-etale site of the indicated
scheme.

Let \(E/\Ql\) be a finite extension contained in \(\Qlb\), let
\(\mathcal O_E\) be its valuation ring, and choose a uniformizer
\(\varpi_E\).  If \(A\) is a finite discrete ring, \(A_X\) denotes the
sheafification of the constant presheaf with value \(A\).  Define
\[
 \widehat{\mathcal O}_{E,X}
   =\varprojlim_n(\mathcal O_E/\varpi_E^n)_X,
 \qquad
 E_X=\widehat{\mathcal O}_{E,X}[1/\varpi_E].
\]

\begin{definition}[Lisse and constructible coefficient sheaves]
An \(E_X\)-module sheaf \(\mathcal F\) is \emph{lisse of rank \(r\)} if
each point of \(X\) has a pro-etale neighbourhood \(U\to X\) for which
\(\mathcal F|_U\simeq E_U^{\oplus r}\).  It is \emph{constructible} if
there is a finite partition
\[
                       X=\coprod_{a=1}^N X_a
\]
by locally closed constructible subschemes such that
\(\mathcal F|_{X_a}\) is lisse of finite rank.  Write
\(\mathsf{Cons}(X,E)\) for the resulting abelian category.
\end{definition}

Put \(A_n=\mathcal O_E/\varpi_E^n\) and let \((A_n)_X\) be the associated
finite constant coefficient sheaf.

\begin{definition}[Finite-torsion constructible categories]
For \(n\ge1\), an \((A_n)_X\)-module sheaf is \emph{locally constant
finite} if every point has a pro-etale neighbourhood on which it is the
constant sheaf associated with a finite, equivalently finitely generated,
\(A_n\)-module.  It is \emph{lisse free of rank \(r\)} if that module may
be chosen to be \(A_n^r\).  A sheaf is constructible if its restriction to
every member of a finite locally closed constructible partition of \(X\)
is locally constant finite.  Write \(\mathsf{Cons}(X,A_n)\) for this
abelian category; kernels and cokernels need not be lisse free, but remain
locally constant finite after refinement.  Let \(D(X,A_n)\) be the ordinary
derived category of all \((A_n)_X\)-module sheaves and define
\[
 D^b_c(X,A_n)=
 \{K\in D(X,A_n):\mathcal H^i(K)\in\mathsf{Cons}(X,A_n),
          \ \mathcal H^i(K)=0\text{ for }|i|\gg0\}.
\]
This is a full ordinary triangulated subcategory of \(D(X,A_n)\).
\end{definition}

\begin{definition}[Bounded constructible derived category]
Let \(D(X,E)\) be the ordinary derived category of sheaves of
\(E_X\)-modules, obtained by localizing cochain complexes at
quasi-isomorphisms.  Define
\[
 \Dbc(X,E)=
 \{K\in D(X,E):\mathcal H^i(K)\in\mathsf{Cons}(X,E)
       \text{ for all }i,\ \mathcal H^i(K)=0\text{ for }|i|\gg0\}.
\]
It is a full triangulated subcategory of \(D(X,E)\).
\end{definition}

\begin{definition}[The abelian category of reduction towers]
Let \(\mathsf{Tow}_E(X)\) be the abelian category whose objects are
inverse systems \(M_\bullet=(M_n,t_{n+1,n})_{n\ge1}\), where \(M_n\) is a
sheaf of \((A_n)_X\)-modules and the \(\mathcal O_E\)-linear transition
map \(t_{n+1,n}:M_{n+1}\to M_n\) is compatible with
\(A_{n+1}\twoheadrightarrow A_n\).  Kernels and cokernels are formed
levelwise.  Write \(D(\mathsf{Tow}_E(X))\) for the ordinary derived
category of this abelian category and
\[
             \operatorname{ev}_n:D(\mathsf{Tow}_E(X))
                    \longrightarrow D(X,A_n)
\]
for derived evaluation at level \(n\).
\end{definition}

\begin{definition}[Normalized bounded adic category]
An object \(K_\bullet\in D(\mathsf{Tow}_E(X))\) is normalized, bounded,
and constructible if \(K_n=\operatorname{ev}_n(K_\bullet)\) belongs to
\(D^b_c(X,A_n)\), one interval \([a,b]\subset\mathbf Z\) contains the
amplitude of every \(K_n\), and the transition induces an isomorphism
\[
 \theta_n:K_{n+1}\otimes^L_{A_{n+1}}A_n
                         \xrightarrow{\sim}K_n
 \qquad(n\ge1).
\]
The two reductions from level \(n+2\) to level \(n\) must agree after the
canonical associativity isomorphism for derived tensor product.  Define
\[
 D^b_c(X,\mathcal O_E)_{\mathrm{adic}}
\]
to be the full subcategory of \(D(\mathsf{Tow}_E(X))\) on these objects.
\end{definition}

\begin{theorem}[Triangulation of normalized towers]
The normalized bounded constructible subcategory is closed under shifts,
isomorphisms, and the third object of a distinguished triangle in
\(D(\mathsf{Tow}_E(X))\).  With the induced distinguished triangles it is
an ordinary triangulated category.  Every evaluation functor
\[
 D^b_c(X,\mathcal O_E)_{\mathrm{adic}}\longrightarrow D^b_c(X,A_n)
\]
is exact.
\end{theorem}

\begin{definition}[Derived inverse limit and derived completion]
Let \(L\in D(X,\widehat{\mathcal O}_{E,X})\), and put
\(L_n=L\otimes^L_{\widehat{\mathcal O}_{E,X}}(A_n)_X\).  Choose a
K-injective strict tower \(I_{n+1}\to I_n\) representing the reduction
tower and define
\[
 \mathbf R\!\varprojlim_nL_n
 =\operatorname{Cone}\!\left(
  1-s:\prod_{n\ge1}I_n\longrightarrow\prod_{n\ge1}I_n
 \right)[-1],
\]
where \(s((x_n)_n)=(I_{n+1}\to I_n)(x_{n+1})\).  The compatible quotient
maps define the completion unit
\[
                   \eta_L:L\longrightarrow\mathbf R\!\varprojlim_nL_n.
\]
The object \(L\) is \emph{derived complete} if \(\eta_L\) is an
isomorphism.  It is \emph{bounded constructibly derived complete} if it is
derived complete, the reductions \(L_n\) are constructible, their
amplitudes lie in one finite interval, and their transition maps are
normalized.  Write
\[
 D^b_{c,\mathrm{dc}}(X,\widehat{\mathcal O}_E)
\]
for the full ordinary triangulated subcategory of such objects.
\end{definition}

\begin{theorem}[Adic realization]
For a normalized tower \(K_\bullet\), regard every \(K_n\) as a complex of
\(\widehat{\mathcal O}_{E,X}\)-modules and define
\[
                \operatorname{Real}(K_\bullet)
                    =\mathbf R\!\varprojlim_nK_n.
\]
Reduction of \(\operatorname{Real}(K_\bullet)\) to \(A_n\) is canonically
isomorphic to \(K_n\).  Realization and reduction are mutually
quasi-inverse exact equivalences
\[
 D^b_c(X,\mathcal O_E)_{\mathrm{adic}}
   \xrightarrow{\sim}
 D^b_{c,\mathrm{dc}}(X,\widehat{\mathcal O}_E).
\]
The equivalence is independent of the selected K-injective strict tower.
\end{theorem}

\begin{definition}[Rationalization, rational category, and lattice]
For \(L\in D^b_{c,\mathrm{dc}}(X,\widehat{\mathcal O}_E)\), define
\[
 L[1/\varpi_E]
   =L\otimes^L_{\widehat{\mathcal O}_{E,X}}E_X
   \simeq\varinjlim\bigl(L\xrightarrow{\varpi_E}L
                  \xrightarrow{\varpi_E}L\longrightarrow\cdots\bigr).
\]
Localize the morphism groups of the derived-complete category by
\(\varpi_E\), and then take the Karoubi envelope.  Via adic realization,
denote the resulting ordinary idempotent-complete triangulated category by
\[
 D^b_c(X,\mathcal O_E)_{\mathrm{adic}}[1/\varpi_E]^\natural.
\]
An \emph{integral lattice} in \(K\in\Dbc(X,E)\) is an object
\(L\in D^b_{c,\mathrm{dc}}(X,\widehat{\mathcal O}_E)\) together with an
isomorphism \(L[1/\varpi_E]\xrightarrow{\sim}K\) in
\(D(X,E)\).
\end{definition}

\begin{theorem}[Coefficient realization]
Every object of \(\Dbc(X,E)\) admits an integral lattice.  Reduction of a
lattice defines exact functors
\[
 D^b_c(X,\mathcal O_E)_{\mathrm{adic}}
   \longrightarrow D^b_c(X,A_n),
 \qquad K_\bullet\longmapsto\operatorname{ev}_n(K_\bullet).
\]
Under adic realization this value is canonically
\(\operatorname{Real}(K_\bullet)\otimes^L_{\widehat{\mathcal O}_{E,X}}
(A_n)_X\).
Rationalization followed by Karoubi completion gives an exact equivalence
\[
 D^b_c(X,\mathcal O_E)_{\mathrm{adic}}[1/\varpi_E]^\natural
        \xrightarrow{\sim}\Dbc(X,E).
\]
Finite scalar extension \(E\subset M\subset\Qlb\) defines an exact functor
\[
             -\otimes_E M:\Dbc(X,E)\longrightarrow\Dbc(X,M).
\]
\end{theorem}

Fix an abstract field embedding
\[
                         \iota:\Qlb\hookrightarrow\mathbf C.
\]
No compatibility with the topologies of the two fields is imposed.

\section{The three Frobenius conventions}

\begin{definition}[Arithmetic and geometric field Frobenius]
The arithmetic Frobenius of \(\bar k/k\) is
\[
 \Fr_{\mathrm{ar}}:\bar k\longrightarrow\bar k,
 \qquad a\longmapsto a^q.
\]
The geometric Frobenius is its inverse
\[
                 \Fr_{\mathrm{geom}}=\Fr_{\mathrm{ar}}^{-1}.
\]
For \(e\ge1\), local arithmetic Frobenius over \(\mathbf F_{q^e}\) is
\(a\mapsto a^{q^e}\), and local geometric Frobenius is its inverse.
\end{definition}

\begin{definition}[Scheme Frobenius]
Form the Frobenius twist
\[
              X^{(q)}=X\times_{\bar k,\Fr_{\mathrm{ar}}}\bar k.
\]
The relative \(q\)-power Frobenius is a morphism
\(F^{(q)}_{X/\bar k}:X\to X^{(q)}\).  The descent of \(X\) from \(X_0\)
gives a canonical isomorphism
\(\delta_X:X^{(q)}\xrightarrow{\sim}X\).  Define
\[
       \Phi_X=\delta_X\circ F^{(q)}_{X/\bar k}:X\longrightarrow X.
\]
This is a universal homeomorphism.  It is generally not an automorphism of
schemes, even though its action on the pro-etale constructible category is
an equivalence.
\end{definition}

\begin{theorem}[Naturality, iteration, and topological invariance]
For every morphism \(f:X_0\to Y_0\), with geometric base change again
denoted \(f:X\to Y\), one has
\[
                         f\circ\Phi_X=\Phi_Y\circ f.
\]
Write \(\Phi_X^{[0]}=\operatorname{id}_X\) and
\(\Phi_X^{[n+1]}=\Phi_X\circ\Phi_X^{[n]}\).  Then
\[
 \Phi_X^{[m+n]}=\Phi_X^{[m]}\circ\Phi_X^{[n]}.
\]
Pullback by \(\Phi_X\) is an exact equivalence on constructible sheaves and
an exact autoequivalence
\[
                  \Phi_X^*:\Dbc(X,E)\xrightarrow{\sim}\Dbc(X,E).
\]
The fixed geometric points of \(\Phi_X^{[n]}\) are canonically
\(X_0(\mathbf F_{q^n})\).  Under the usual descent comparison for sheaves,
\(\Phi_X^*\) corresponds to geometric Galois Frobenius.
\end{theorem}

\begin{definition}[Ordered iterate of a Frobenius structure]
Let \(\mathcal F\in\mathsf{Cons}(X,E)\), and let
\(\varphi:\Phi_X^*\mathcal F\xrightarrow{\sim}\mathcal F\).  Define
\[
 \varphi^{[0]}=\operatorname{id}_{\mathcal F},\qquad
 \varphi^{[n+1]}
  =\varphi^{[n]}\circ(\Phi_X^{[n]})^*\varphi:
  (\Phi_X^{[n+1]})^*\mathcal F\longrightarrow\mathcal F,
\]
using the canonical comparison between pullback by a composite and the
composite of pullbacks.  These iterates satisfy
\[
 \varphi^{[m+n]}
 =\varphi^{[m]}\circ(\Phi_X^{[m]})^*\varphi^{[n]}.
\]
The same definition applies degree by degree to a strict complex.
\end{definition}

\begin{definition}[Closed-point geometric Frobenius]
Let \(x\in|X_0|\) be a closed point.  Put
\[
 d(x)=[\kappa(x):k],\qquad q_x=\#\kappa(x)=q^{d(x)}.
\]
Choose a geometric point \(\bar x\to x\).  Its orbit under \(\Phi_X\) has
length \(d(x)\), and \(\Phi_X^{[d(x)]}(\bar x)=\bar x\).  The stalk of
\(\varphi^{[d(x)]}\) at \(\bar x\), after the canonical stalk-pullback
identification, is an automorphism
\[
              F_x:\mathcal F_{\bar x}\longrightarrow\mathcal F_{\bar x}.
\]
It is local geometric Frobenius.  Replacing \(\bar x\) conjugates \(F_x\),
so its characteristic polynomial and eigenvalue multiset are canonical.
\end{definition}

\begin{proposition}[Local and global iterate comparison]
Let \(n\ge1\), let \(y\in X_0(\mathbf F_{q^n})\), and let \(x\) be the
closed point below \(y\).  Then \(d(x)\mid n\).  On the stalk at \(y\), the
operator induced by \(\varphi^{[n]}\) is conjugate to
\[
                              F_x^{n/d(x)}.
\]
Thus the fixed-point operator and the closed-point operator use the same
geometric Frobenius normalization.
\end{proposition}

\section{Strict Weil categories}

\begin{definition}[Constructible Weil sheaf]
A constructible \(E\)-Weil sheaf on \(X_0\) is a pair
\((\mathcal F,\varphi_{\mathcal F})\) with
\(\mathcal F\in\mathsf{Cons}(X,E)\) and an isomorphism
\[
       \varphi_{\mathcal F}:\Phi_X^*\mathcal F\xrightarrow{\sim}\mathcal F.
\]
A morphism \(u:(\mathcal F,\varphi_{\mathcal F})\to
(\mathcal G,\varphi_{\mathcal G})\) is a sheaf morphism satisfying
\[
 u\circ\varphi_{\mathcal F}
  =\varphi_{\mathcal G}\circ\Phi_X^*u.
\]
The resulting abelian category is denoted
\(\mathsf{Cons}_{\mathrm W}(X_0,E)\).  Kernels, images, and cokernels are
taken on the underlying sheaves and carry their induced Weil structures.
\end{definition}

\begin{definition}[Strict bounded Weil complex]
A strict bounded Weil complex is a bounded cochain complex in
\(\mathsf{Cons}_{\mathrm W}(X_0,E)\).  Explicitly it consists of a bounded
complex \((K^i,d^i)\), isomorphisms
\(\varphi^i:\Phi_X^*K^i\to K^i\), and the equations
\[
          d^i\circ\varphi^i=\varphi^{i+1}\circ\Phi_X^*(d^i).
\]
A strict morphism commutes with the Frobenius isomorphisms degree by degree.
It is an equivariant quasi-isomorphism if its underlying map of sheaf
complexes induces an isomorphism on every cohomology sheaf.
\end{definition}

\begin{definition}[Finite-stage derived Weil category]
Let \(K^b(\mathcal A)\) be the homotopy category of bounded complexes in an
abelian category \(\mathcal A\).  Define
\[
 \WeilD(X_0,E)
 =K^b\bigl(\mathsf{Cons}_{\mathrm W}(X_0,E)\bigr)
   [\text{equivariant quasi-isomorphisms}^{-1}].
\]
This is the ordinary Verdier localization of strict equivariant complexes.
It is not defined as the category of pairs
\((K,\Phi_X^*K\simeq K)\) formed after passing to \(\Dbc(X,E)\).
\end{definition}

\begin{theorem}[Triangulation and cohomology]
The category \(\WeilD(X_0,E)\) is an ordinary triangulated category.  Its
shift and cone are induced by the shift and cone of strict complexes.  The
forgetful functor
\[
 \operatorname{forg}_E:\WeilD(X_0,E)\longrightarrow\Dbc(X,E)
\]
is exact.  For every \(K\), the sheaf \(\mathcal H^i(K)\) carries the Weil
structure induced from a strict representative, independently of that
representative.
\end{theorem}

\begin{theorem}[Finite scalar extension]
For \(E\subset M\subset\Qlb\) finite over \(\Ql\), degreewise scalar
extension and the isomorphism
\[
 \Phi_X^*(K\otimes_E M)\simeq\Phi_X^*K\otimes_E M
\]
define an exact functor
\[
 -\otimes_E M:\WeilD(X_0,E)\longrightarrow\WeilD(X_0,M).
\]
It commutes with shifts, cones, cohomology, stalks, and closed-point
Frobenius operators.
\end{theorem}

\begin{definition}[Total complexes]
For a double complex \(B^{p,q}\) with commuting raw differentials
\(d_h:B^{p,q}\to B^{p+1,q}\) and
\(d_v:B^{p,q}\to B^{p,q+1}\), require
\[
 d_h^2=d_v^2=0,\qquad d_hd_v=d_vd_h,
\]
and define the direct-sum total complex by
\[
 (\operatorname{Tot}^{\oplus}B)^n=\bigoplus_{p+q=n}B^{p,q},
 \qquad d|_{B^{p,q}}=d_h+(-1)^p d_v.
\]
Define the product total complex by replacing the direct sum with
\(\prod_{p+q=n}\).  For bounded bicomplexes the two definitions agree.
For complexes \(P,I\), the internal-Hom complex is the product total
complex
\[
 \mathcal Hom^n(P,I)=\prod_p\mathcal Hom(P^p,I^{p+n}),
 \qquad d(f)=d_I f-(-1)^n f d_P.
\]
\end{definition}

\begin{definition}[Ambient equivariant derived category]
Let \(\mathsf{Mod}_{\mathrm W}(X_0,E)\) be the abelian category of all
\(E_X\)-module sheaves \(M\), without a constructibility hypothesis,
equipped with an isomorphism
\(\varphi_M:\Phi_X^*M\xrightarrow{\sim}M\).  Its morphisms commute with
these isomorphisms.  Let
\[
 D_{\mathrm W}(X_0,E)=D\bigl(\mathsf{Mod}_{\mathrm W}(X_0,E)\bigr)
\]
be the ordinary unbounded derived category.  Denote by
\(D^b_{\mathrm W,c,\mathrm{amb}}(X_0,E)\) the full triangulated subcategory
whose underlying cohomology sheaves are constructible and vanish outside a
finite interval.  Inclusion of constructible equivariant sheaves induces
\[
 \operatorname{Real}_{\mathrm W}:\WeilD(X_0,E)
       \longrightarrow D^b_{\mathrm W,c,\mathrm{amb}}(X_0,E).
\]
Write
\[
 U:\mathsf{Mod}_{\mathrm W}(X_0,E)\longrightarrow
       \mathsf{Mod}(X,E)
\]
for the functor which forgets the Weil isomorphism.
\end{definition}

\begin{theorem}[Grothendieck property of equivariant sheaves]
The category \(\mathsf{Mod}_{\mathrm W}(X_0,E)\) is Grothendieck abelian.
The functor \(U\) is exact and conservative, commutes with all small
colimits, and sends injective objects to objects which are acyclic for
\(\mathcal Hom_{E_X}(A,-)\) whenever \(A\) is a bounded complex of
constructible finite-rank coefficient sheaves.
\end{theorem}

\begin{definition}[Equivariant adapted replacement classes]
Let \(C_{\mathrm W}(X_0,E)\) be the category of unbounded complexes in
\(\mathsf{Mod}_{\mathrm W}(X_0,E)\).  An equivariant quasi-isomorphism
\(q:P\to K\) belongs to \(\mathcal Q_{\mathrm W}(K)\) if \(U(P)\) is
K-flat, meaning that
\[
 \operatorname{Tot}^{\oplus}(U(P)\otimes_{E_X}A)
 \quad\hbox{is acyclic for every acyclic complex }A.
\]
An equivariant quasi-isomorphism \(\eta:K\to I\) belongs to
\(\mathcal I_{\mathrm W}(K)\) if \(U(I)\) is K-injective, meaning that
\[
 \mathcal Hom_{E_X}^{\bullet}(A,U(I))
 \quad\hbox{is acyclic for every acyclic complex }A.
\]
A morphism in either replacement class is a chain map over, respectively
under, \(K\).  The same definitions, with the Weil structures forgotten,
give the classes \(\mathcal Q(K)\) and \(\mathcal I(K)\) in
\(C(\mathsf{Mod}(X,E))\).
\end{definition}

\begin{theorem}[Functorial equivariant adapted replacements]
There are functors
\[
 Q_{\mathrm W},J_{\mathrm W}:C_{\mathrm W}(X_0,E)
                  \longrightarrow C_{\mathrm W}(X_0,E)
\]
and natural transformations
\[
 q:Q_{\mathrm W}\Longrightarrow\operatorname{id},
 \qquad \eta:\operatorname{id}\Longrightarrow J_{\mathrm W}
\]
such that \(q_K\in\mathcal Q_{\mathrm W}(K)\) and
\(\eta_K\in\mathcal I_{\mathrm W}(K)\) for every \(K\).  Both functors
send equivariant quasi-isomorphisms to equivariant quasi-isomorphisms.
For every equivariant quasi-isomorphism \(u\), the map
\(J_{\mathrm W}(u)\) is a homotopy equivalence between K-injective
complexes.
For every finite extension \(E\subset M\), there are natural
quasi-isomorphisms
\[
 Q_{\mathrm W,E}(K)\otimes_E M\longrightarrow
 Q_{\mathrm W,M}(K\otimes_E M),
 \qquad
 J_{\mathrm W,E}(K)\otimes_E M\longrightarrow
 J_{\mathrm W,M}(K\otimes_E M),
\]
compatible with \(q,\eta\), identities, and towers of coefficient fields.
\end{theorem}

\begin{definition}[Equivariant Ext comparison]
For \(\mathcal F,\mathcal G\in\mathsf{Cons}_{\mathrm W}(X_0,E)\), let
\(\operatorname{Ext}^r_{\mathrm W,c}(\mathcal F,\mathcal G)\) be Yoneda
Ext in \(\mathsf{Cons}_{\mathrm W}(X_0,E)\), and let
\(\operatorname{Ext}^r_{\mathrm W,amb}(\mathcal F,\mathcal G)\) be Yoneda
Ext in \(\mathsf{Mod}_{\mathrm W}(X_0,E)\).  The exact inclusion defines
\[
 e^r_{\mathcal F,\mathcal G}:
 \operatorname{Ext}^r_{\mathrm W,c}(\mathcal F,\mathcal G)
 \longrightarrow
 \operatorname{Ext}^r_{\mathrm W,amb}(\mathcal F,\mathcal G).
\]
\end{definition}

\begin{theorem}[Equivariant Ext effacement]
The category \(\mathsf{Cons}_{\mathrm W}(X_0,E)\) is a Serre subcategory of
\(\mathsf{Mod}_{\mathrm W}(X_0,E)\).  For every
\(\mathcal F,\mathcal G\in\mathsf{Cons}_{\mathrm W}(X_0,E)\) and every
\(r>0\), every class
\(\xi\in\operatorname{Ext}^r_{\mathrm W,amb}(\mathcal F,\mathcal G)\)
admits a constructible equivariant monomorphism
\(\mathcal G\hookrightarrow\mathcal G'\) for which
\[
 \operatorname{Ext}^r_{\mathrm W,amb}(\mathcal F,\mathcal G)
 \longrightarrow
 \operatorname{Ext}^r_{\mathrm W,amb}(\mathcal F,\mathcal G')
\]
sends \(\xi\) to zero.  Consequently
\(e^r_{\mathcal F,\mathcal G}\) is an isomorphism
for every \(r\ge0\), naturally in both variables and compatibly with Yoneda
products and finite scalar extension.
\end{theorem}

\begin{theorem}[Full faithfulness of bounded realization]
For bounded strict constructible Weil complexes \(K,L\), the localization
map is a bijection
\[
 \operatorname{Hom}_{\WeilD(X_0,E)}(K,L)
 \xrightarrow{\sim}
 \operatorname{Hom}_{D_{\mathrm W}(X_0,E)}
       (\operatorname{Real}_{\mathrm W}K,
        \operatorname{Real}_{\mathrm W}L).
\]
The same assertion holds after replacing \(L\) by any shift \(L[r]\), and
these bijections respect composition, shifts, and connecting morphisms.
\end{theorem}

\begin{theorem}[Bounded strictification by cohomological amplitude]
Let \(M\in D_{\mathrm W}(X_0,E)\) have constructible cohomology and suppose
\(\mathcal H^i(M)=0\) outside \([a,b]\).  There is a bounded strict
constructible Weil complex \(K_M\), obtained by finitely many extensions of
shifts of constructible Weil sheaves and having cohomology concentrated in
\([a,b]\), together with an isomorphism
\[
 \operatorname{Real}_{\mathrm W}(K_M)\xrightarrow{\sim}M.
\]
For a morphism \(M\to N\), the full-faithfulness theorem gives a unique
morphism \(K_M\to K_N\) after the selected realization isomorphisms are
fixed.  These assignments preserve identities and composition.
\end{theorem}

\begin{theorem}[Bounded strict equivariant realization]
The functor \(\operatorname{Real}_{\mathrm W}\) is an equivalence of
ordinary triangulated categories.  It admits a triangulated quasi-inverse
\[
 \operatorname{Str}_{\mathrm W}:
 D^b_{\mathrm W,c,\mathrm{amb}}(X_0,E)\longrightarrow\WeilD(X_0,E)
\]
with natural unit and counit satisfying the two triangle identities.  The
equivalence and the selected unit and counit commute with the forgetful
functor, finite scalar extension, shifts, and distinguished
triangles, and cohomology sheaves.
\end{theorem}

\begin{definition}[Derived tensor product and internal Hom]
For \(K,L\in D(X,E)\), choose K-flat quasi-isomorphisms \(P_K\to K\) and
\(P_L\to L\), and choose a K-injective quasi-isomorphism \(L\to I_L\).
Here K-flat means that \(P\otimes_{E_X}A\) is acyclic for every acyclic
complex \(A\), and K-injective means that
\(\mathcal Hom_{E_X}(A,I)\) is acyclic for every acyclic complex \(A\).
Define
\[
 K\otimes_E^LL=\operatorname{Tot}^{\oplus}(P_K\otimes_{E_X}P_L),
 \qquad
 \RHom_E(K,L)=\mathcal Hom^\bullet_{E_X}(P_K,I_L).
\]
The total differential on homogeneous tensors is
\(d(x\otimes y)=d(x)\otimes y+(-1)^{\deg x}x\otimes d(y)\), and the
internal-Hom differential is
\(d(f)=d_{I_L}f-(-1)^{\deg f}f d_{P_K}\).  Independence of the selected
replacements gives bifunctors
\[
 \begin{aligned}
 -\otimes_E^L- &:D(X,E)\times D(X,E)\longrightarrow D(X,E),\\
 \RHom_E(-,-) &:D(X,E)^{\mathrm{op}}\times D(X,E)
                                      \longrightarrow D(X,E).
 \end{aligned}
\]
\end{definition}

\begin{theorem}[Finite Tor amplitude and bounded constructibility]
Every object of \(\Dbc(X,E)\) is locally represented on a finite
constructible stratification by a bounded complex of finite-rank locally
free \(E_X\)-modules.  It therefore has finite Tor amplitude.  Consequently
the preceding bifunctors restrict to
\[
 \begin{aligned}
 -\otimes_E^L- &: \Dbc(X,E)\times\Dbc(X,E)\longrightarrow\Dbc(X,E),\\
 \RHom_E(-,-) &: \Dbc(X,E)^{\mathrm{op}}\times\Dbc(X,E)
                                      \longrightarrow\Dbc(X,E).
 \end{aligned}
\]
They are exact in each variable and commute with finite scalar extension.
\end{theorem}

\begin{proposition}[Closed monoidal Frobenius pullback]
The exact autoequivalence \(\Phi_X^*\) is strong symmetric monoidal and
closed.  It has coherent natural isomorphisms
\[
 c^\otimes_{M,N}:\Phi_X^*(M\otimes_{E_X}N)
   \xrightarrow{\sim}\Phi_X^*M\otimes_{E_X}\Phi_X^*N,
\]
\[
 c^{\mathcal Hom}_{M,N}:
 \Phi_X^*\mathcal Hom_{E_X}(M,N)
 \xrightarrow{\sim}
 \mathcal Hom_{E_X}(\Phi_X^*M,\Phi_X^*N).
\]
They commute with the differentials of tensor and internal-Hom complexes
and are compatible with associativity, symmetry, the coefficient unit,
evaluation, currying, and finite scalar extension.
\end{proposition}

\begin{definition}[Ambient equivariant derived tensor and internal Hom]
For complexes \(K,L\in C_{\mathrm W}(X_0,E)\), define
\[
 K\otimes_{\mathrm W,E}^LL
 =\operatorname{Tot}^{\oplus}
    \bigl(UQ_{\mathrm W}K\otimes_{E_X}UQ_{\mathrm W}L\bigr),
\]
\[
 R\mathcal Hom_{\mathrm W,E}(K,L)
 =\mathcal Hom_{E_X}^\bullet
    \bigl(UQ_{\mathrm W}K,UJ_{\mathrm W}L\bigr).
\]
The first complex has the Weil arrow
\[
 \Phi_X^*\operatorname{Tot}^{\oplus}
   (UQ_{\mathrm W}K\otimes UQ_{\mathrm W}L)
 \xrightarrow{c^\otimes}
 \operatorname{Tot}^{\oplus}
   (\Phi_X^*UQ_{\mathrm W}K\otimes\Phi_X^*UQ_{\mathrm W}L)
 \xrightarrow{\operatorname{Tot}(\varphi_{QK}\otimes\varphi_{QL})}
 \operatorname{Tot}^{\oplus}
   (UQ_{\mathrm W}K\otimes UQ_{\mathrm W}L),
\]
Put
\[
 \psi_{K,L}=\mathcal Hom^\bullet
   (\varphi_{QK}^{-1},\varphi_{JL}).
\]
The second complex has the Weil arrow
\[
 \begin{aligned}
 &\Phi_X^*\mathcal Hom_{E_X}^\bullet
   (UQ_{\mathrm W}K,UJ_{\mathrm W}L)\\
 &\quad\xrightarrow{c^{\mathcal Hom}}
 \mathcal Hom_{E_X}^\bullet
   (\Phi_X^*UQ_{\mathrm W}K,\Phi_X^*UJ_{\mathrm W}L)\\
 &\quad\xrightarrow{\psi_{K,L}}
 \mathcal Hom_{E_X}^\bullet
   (UQ_{\mathrm W}K,UJ_{\mathrm W}L).
 \end{aligned}
\]
Here \(\varphi_{QK},\varphi_{QL},\varphi_{JL}\) are the degreewise descent
isomorphisms belonging to the actual equivariant replacement complexes.
\end{definition}

\begin{theorem}[Descent of the ambient equivariant operations]
The four maps
\[
 Q_{\mathrm W}(u)\otimes Q_{\mathrm W}(L),\quad
 Q_{\mathrm W}(K)\otimes Q_{\mathrm W}(v),\quad
 \mathcal Hom^\bullet(Q_{\mathrm W}(u),J_{\mathrm W}(L)),\quad
 \mathcal Hom^\bullet(Q_{\mathrm W}(K),J_{\mathrm W}(v))
\]
are quasi-isomorphisms whenever \(u\) or \(v\), in its displayed
contravariant or covariant position, is an equivariant quasi-isomorphism.
Consequently the preceding formulas define bifunctors on
\(D_{\mathrm W}(X_0,E)\).  After applying \(U\), they are canonically the
ordinary derived tensor product and derived internal Hom.
\end{theorem}

\begin{theorem}[Equivariant bounded constructible closure]
For \(K,L\in D^b_{\mathrm W,c,\mathrm{amb}}(X_0,E)\), both
\(K\otimes_{\mathrm W,E}^LL\) and
\(R\mathcal Hom_{\mathrm W,E}(K,L)\) again belong to
\(D^b_{\mathrm W,c,\mathrm{amb}}(X_0,E)\).  This restriction is exact in
each variable and commutes with finite scalar extension.
\end{theorem}

\begin{definition}[Transported strict tensor and internal Hom]
For \(K,L\in\WeilD(X_0,E)\), define
\[
 K\otimes_E^LL=\operatorname{Str}_{\mathrm W}
  \bigl(\operatorname{Real}_{\mathrm W}K
       \otimes_{\mathrm W,E}^L\operatorname{Real}_{\mathrm W}L\bigr),
\]
\[
 \RHom_E(K,L)=\operatorname{Str}_{\mathrm W}
  R\mathcal Hom_{\mathrm W,E}
  \bigl(\operatorname{Real}_{\mathrm W}K,
        \operatorname{Real}_{\mathrm W}L\bigr).
\]
Their strict Weil arrows are the images under
\(\operatorname{Str}_{\mathrm W}\) of the two typed arrows in the preceding
definition.
\end{definition}

\begin{theorem}[Strict tensor, internal Hom, and coherence]
The transported operations are exact bifunctors
\[
\begin{aligned}
 -\otimes_E^L- &: \WeilD(X_0,E)\times\WeilD(X_0,E)
                                   \longrightarrow\WeilD(X_0,E),\\
 \RHom_E(-,-) &: \WeilD(X_0,E)^{\mathrm{op}}\times\WeilD(X_0,E)
                                   \longrightarrow\WeilD(X_0,E).
\end{aligned}
\]
Their associator, symmetry, unit, evaluation, tensor--Hom adjunction, shift,
distinguished-triangle, and finite-scalar-extension isomorphisms are transported through
\(\operatorname{Real}_{\mathrm W}\).  They satisfy the pentagon, hexagon,
triangle, adjunction triangle, identity, and scalar-transitivity equations.
Their forgetful images are the corresponding bounded constructible derived
operations.
Under the realization equivalence, their Weil arrows identify with
\[
 \Phi_X^*(K\otimes_E^LL)\xrightarrow{c^\otimes}
 \Phi_X^*K\otimes_E^L\Phi_X^*L
 \xrightarrow{\varphi_K\otimes\varphi_L}K\otimes_E^LL
\]
and
\[
 \Phi_X^*\RHom_E(K,L)\xrightarrow{c^{\mathcal Hom}}
 \RHom_E(\Phi_X^*K,\Phi_X^*L)
 \xrightarrow{\RHom_E(\varphi_K^{-1},\varphi_L)}\RHom_E(K,L).
\]
\end{theorem}

\begin{definition}[Strict linear dual]
The coefficient unit \(E_X\) has the strict identity descent arrow
\(\varphi_E:\Phi_X^*E_X\xrightarrow{\sim}E_X\).  Define
\[
                         K^\vee=\RHom_E(K,E_X).
\]
Its strict Weil arrow is
\[
 \Phi_X^*K^\vee
 \xrightarrow{\sim}\RHom_E(\Phi_X^*K,\Phi_X^*E_X)
 \xrightarrow{\RHom_E(\varphi_K^{-1},\varphi_E)}K^\vee.
\]
For a lisse finite-rank sheaf this agrees with the usual sheaf dual and
inverse-transpose Frobenius.
\end{definition}

\begin{definition}[The finite-stage \(\Qlb\)-category]
First form an additive category whose objects are pairs \((E,K)\), where
\(E/\Ql\) is finite inside \(\Qlb\) and \(K\in\WeilD(X_0,E)\), and whose
morphism groups are
\[
 \operatorname{Hom}((E,K),(E',L))
 =\varinjlim_{M\supset E,E'}
 \operatorname{Hom}_{\WeilD(X_0,M)}
     (K\otimes_E M,L\otimes_{E'}M).
\]
Composition, shifts, cones, and distinguished triangles are formed after
passing all data to one common finite coefficient extension.  Take the
Karoubi envelope: its objects are pairs \((P,e)\) with \(e^2=e\), and its
morphisms \(u:(P,e)\to(Q,d)\) satisfy \(u=due\).  Denote the resulting
ordinary idempotent-complete triangulated category by
\[
                  \WeilD(X_0,\Qlb)_{\mathrm{2-colim}}.
\]
A triangle in the Karoubi envelope is distinguished if it is isomorphic to
a direct summand of a distinguished triangle represented at one finite
stage.
\end{definition}

\begin{theorem}[Finite-stage descent]
Every finite collection of objects, morphisms, distinguished triangles,
idempotents, cohomology objects, tensor products, internal Homs, and linear
duals in
\(\WeilD(X_0,\Qlb)_{\mathrm{2-colim}}\) is represented over a single finite
extension \(E/\Ql\).  Two such representations agree exactly when they
agree after extension to a common finite coefficient field \(M\).  For a retract
\((E,K,e)\), its \(i\)-th cohomology object is the image of
\(\mathcal H^i(e)\) in the Karoubi envelope of constructible Weil sheaves.
\end{theorem}

\section{Tate objects and Verdier duality}

\begin{definition}[Tate objects]
On the pro-etale site define
\[
 \mathbf Z_\ell(1)=\varprojlim_n\mu_{\ell^n},
 \qquad
 \mu_{\ell^n}=\ker([\ell^n]:\mathbf G_m\to\mathbf G_m),
\]
where \(\mu_{\ell^{n+1}}\to\mu_{\ell^n}\) is the \(\ell\)-th power map.
Put
\[
 E(1)=\mathbf Z_\ell(1)\otimes_{\mathbf Z_\ell}E,
 \quad E(0)=E_X,
 \quad E(m)=E(1)^{\otimes_E^L m}\ (m>0),
 \quad E(-m)=\RHom_E(E(m),E_X)\ (m>0),
\]
and, for every \(m\in\mathbf Z\), put
\(K(m)=K\otimes_E^L E(m)\).
\end{definition}

\begin{definition}[Strict Weil structure on the Tate object]
Functorial pullback of roots of unity, followed by the descent
identification for scheme Frobenius, gives compatible isomorphisms
\[
 \varphi_{\ell^n}:\Phi_X^*\mu_{\ell^n}
                         \xrightarrow{\sim}\mu_{\ell^n}
 \qquad(n\ge1).
\]
They commute with the transition maps
\(\mu_{\ell^{n+1}}\to\mu_{\ell^n}\), \(\zeta\mapsto\zeta^\ell\).
Define
\[
 \varphi_{\mathbf Z_\ell(1)}
   =\varprojlim_n\varphi_{\ell^n}:
 \Phi_X^*\mathbf Z_\ell(1)\xrightarrow{\sim}\mathbf Z_\ell(1),
\]
and define
\[
 \varphi_{E(1)}=\varphi_{\mathbf Z_\ell(1)}\otimes E:
                  \Phi_X^*E(1)\xrightarrow{\sim}E(1).
\]
Put \(\varphi_{E(0)}=\varphi_E\).  For \(m>0\), define
\(\varphi_{E(m)}\) by the iterated strict tensor arrow, and define the
negative twist by the typed composite
\[
 \varphi_{E(-m)}:\Phi_X^*\RHom_E(E(m),E_X)
 \xrightarrow{c^{\mathcal Hom}}
 \RHom_E(\Phi_X^*E(m),\Phi_X^*E_X)
 \xrightarrow{\RHom_E(\varphi_{E(m)}^{-1},\varphi_E)}
 \RHom_E(E(m),E_X).
\]
For every strict Weil complex \((K,\varphi_K)\), define the arrow on its
twist by
\[
 \varphi_{K(m)}:\Phi_X^*(K\otimes_E^LE(m))
 \xrightarrow{c^\otimes}
 \Phi_X^*K\otimes_E^L\Phi_X^*E(m)
 \xrightarrow{\varphi_K\otimes\varphi_{E(m)}}K\otimes_E^LE(m).
\]
These formulas define objects of \(\WeilD(X_0,E)\) for every integer
\(m\).
\end{definition}

\begin{theorem}[Tate Frobenius normalization]
Let \(x\in|X_0|\).  The ordered
\([\kappa(x):k]\)-fold iterate of \(\varphi_{E(1)}\) induces
\[
               F_x|E(1)_{\bar x}=q_x^{-1}\operatorname{id}.
\]
Consequently \(E(m)\) has local geometric Frobenius
\(q_x^{-m}\operatorname{id}\) for every \(m\in\mathbf Z\), and twisting a
finite Frobenius module multiplies its local Frobenius by \(q_x^{-m}\).
\end{theorem}

\begin{definition}[Finite-coefficient derived tensor and internal Hom]
Let \(A=A_n\).  A complex \(P\) of \(A_X\)-module sheaves is K-flat if
\(\operatorname{Tot}^{\oplus}(P\otimes_{A_X}C)\) is acyclic for every
acyclic complex \(C\); a complex \(I\) is K-injective if
\(\mathcal Hom_{A_X}^\bullet(C,I)\) is acyclic for every acyclic
\(C\).  Choose functorial quasi-isomorphisms
\(P_K\to K\), \(P_L\to L\), and \(L\to I_L\) with these properties and define
\[
 K\otimes_A^LL=\operatorname{Tot}^{\oplus}(P_K\otimes_{A_X}P_L),
 \qquad
 R\mathcal Hom_A(K,L)=\mathcal Hom_{A_X}^\bullet(P_K,I_L).
\]
These formulas descend to exact bifunctors
\[
 \begin{aligned}
 -\otimes_A^L- &:D(X,A)\times D(X,A)\longrightarrow D(X,A),\\
 R\mathcal Hom_A(-,-)&:D(X,A)^{\mathrm{op}}\times D(X,A)
                          \longrightarrow D(X,A).
 \end{aligned}
\]
The evaluation, currying, and tensor--Hom adjunction maps are those induced
by the corresponding maps of the selected complexes.
\end{definition}

\begin{definition}[Finite-Tor constructible category]
For a coefficient ring \(A\), a complex \(K\in D(X,A)\) has Tor amplitude
in \([r,s]\) if, for every sheaf \(M\) of \(A_X\)-modules,
\[
 \mathcal H^i(K\otimes_A^LM)=0\qquad(i\notin[r,s]).
\]
Let \(D_{\mathrm{ctf}}(X,A)\) be the full subcategory of
\(D^b_c(X,A)\) consisting of finite-Tor objects.  Since \(E\) is a field,
\[
                 D_{\mathrm{ctf}}(X,E)=\Dbc(X,E).
\]
\end{definition}

\begin{theorem}[Finite-coefficient perfect closure]
The full category \(D_{\mathrm{ctf}}(X,A)\) consists locally on a finite
constructible stratification of bounded complexes of finite projective
\(A\)-modules.  It is closed under derived tensor product.  If
\(K,L\in D_{\mathrm{ctf}}(X,A)\), then
\(R\mathcal Hom_A(K,L)\) is bounded constructible and has finite Tor
amplitude.  These operations commute with reduction \(A_n\to A_m\), with
the identity and cocycle compatibilities.
\end{theorem}

\begin{definition}[Structural compact-support functor]
Let \(a_X:X\to S=\Spec\bar k\) be the structural morphism.  Choose a
compactification
\[
 X\xrightarrow{j}\overline X\xrightarrow{\bar a}S
\]
with \(j\) open and \(\bar a\) proper.  For a sheaf \(\mathcal F\) on
\(X\), let \(j_!\mathcal F\) be the sheafification of
\[
 V\longmapsto
 \{s\in\mathcal F(V\times_{\overline X}X):
       \operatorname{supp}(s)\text{ is closed in }V\}.
\]
For any morphism \(g:Z\to Y\), define
\((g_*\mathcal F)(V)=\mathcal F(V\times_YZ)\) and define \(Rg_*\) by a
K-injective resolution.  Put
\[
                  Ra_{X,!}=R\bar a_*\,j_!:D(X,A)\to D(S,A).
\]
\end{definition}

\begin{theorem}[Compactification independence and composition]
The functor \(Ra_{X,!}\) is independent, up to a canonical isomorphism, of
the chosen compactification.  The comparison isomorphisms are the identity
for identical compactifications, are transitive under proper refinements,
and are compatible with shifts and distinguished triangles.  For separated finite-type
morphisms \(X\xrightarrow fY\xrightarrow gZ\), the canonical maps
\[
                 Rg_!Rf_!\xrightarrow{\sim}R(gf)_!
\]
satisfy the identity and associativity equations.
\end{theorem}

\begin{definition}[Full site-local cohomological dimension]
Let \(p:Z\to Y\) be proper, fix a basis \(\mathcal B_Y\) of affine
quasi-compact weakly contractible objects in the selected pro-etale site,
and let \(c\ge0\).  The condition \(\operatorname{FullCD}(p,A;c)\) means
that
\[
 R^rp_*\mathcal F=0\quad(r>c)
\]
for every sheaf \(\mathcal F\) of \(A\)-modules on \(Z\), and that the same
integer \(c\) has this property for \(p_V:Z\times_YV\to V\), every
\(V\in\mathcal B_Y\), and every sheaf on \(Z\times_YV\).
\end{definition}

\begin{theorem}[Uniform full cohomological dimension]
For the proper morphism \(\bar a:\overline X\to S\), there is an integer
\(c\geq0\), depending on \(\bar a\) but independent of \(n\), such that
\(\operatorname{FullCD}(\bar a,A_n;c)\) holds for every \(n\geq1\).
The same bound is valid after pullback to every member of the selected
weakly contractible affine basis.
\end{theorem}

\begin{theorem}[Coproduct preservation and the exceptional adjoint]
For the uniform integer \(c\) in the preceding theorem,
\(Ra_{X,!}\) preserves small coproducts,
and Brown representability gives an exact right adjoint
\[
 Ra_X^!:D(S,A_n)\longrightarrow D(X,A_n),
 \qquad Ra_{X,!}\dashv Ra_X^!,
\]
whose unit and counit satisfy the two triangle identities.
\end{theorem}

\begin{theorem}[Finite-level compact-support amplitude and devissage]
For every \(n\), compactly supported direct image restricts to
\[
 Ra_{X,!}:D_{\mathrm{ctf}}(X,A_n)\longrightarrow
          D_{\mathrm{ctf}}(S,A_n).
\]
There is one finite cohomological-amplitude interval valid for all \(n\),
and the construction commutes with reduction \(A_n\to A_m\).
\end{theorem}

\begin{theorem}[Exceptional constructibility and finite-Tor preservation]
For every finite stage \(A=A_n\), the exceptional adjunction
restricts to exact functors
\[
 Ra_{X,!}:D_{\mathrm{ctf}}(X,A)\rightleftarrows
 D_{\mathrm{ctf}}(S,A):Ra_X^!.
\]
The amplitude of \(Ra_X^!\) is bounded independently of \(n\), and the
adjunction, its unit, and its counit commute with finite reduction.
\end{theorem}

\begin{definition}[Dualizing object and ambient Verdier dual]
For a finite coefficient ring \(A=A_n\), define
\[
 \omega_{X,A}=Ra_X^!A_S,
 \qquad
 D_{X,A}^{\mathrm{amb}}(K)=R\mathcal Hom_A(K,\omega_{X,A}).
\]
The evaluation map and tensor--Hom adjunction define a natural biduality
morphism in the ambient derived category
\[
\delta_{X,A,K}:K\longrightarrow
 D_{X,A}^{\mathrm{amb}}D_{X,A}^{\mathrm{amb}}K.
\]
When no ambiguity is possible, write \(D_{X,A}\) for this ambient functor;
its restriction to finite-Tor constructible objects is established below.
\end{definition}

\begin{definition}[Lisse-perfect complexes and duality stratifications]
A complex \(K\in D(U,A)\) is lisse perfect if, pro-etale locally on \(U\),
it is the constant-sheaf realization of a bounded complex of finite
projective \(A\)-modules.  A duality stratification of
\(K\in D_{\mathrm{ctf}}(X,A)\) is a finite decomposition
\[
 X=\coprod_{r=1}^mX_r
\]
by locally closed subschemes smooth over \(S\), ordered so that
\(\coprod_{r\le t}X_r\) is open for every \(t\), and such that
\(K|_{X_r}\) is lisse perfect.
\end{definition}

\begin{definition}[Finite-stage Tate objects and coefficient dual]
For every \(\bar k\)-scheme \(T\), define the normalized Tate system and
its finite reductions by
\[
 \mathcal O_{E,T}(1)
   =\mathbf Z_{\ell,T}(1)\otimes_{\mathbf Z_\ell}\mathcal O_E,
 \qquad
 A_{n,T}(1)=\mathcal O_{E,T}(1)\otimes_{\mathcal O_E}^LA_n.
\]
Put \(A_{n,T}(0)=A_{n,T}\), define positive twists by derived tensor
powers, and define
\[
 A_{n,T}(-d)=R\mathcal Hom_{A_n}(A_{n,T}(d),A_{n,T})
 \qquad(d>0).
\]
For \(K\in D(T,A_n)\), define its coefficient dual by
\[
                         K^{\vee_{A_n}}
 =R\mathcal Hom_{A_n}(K,A_{n,T}).
\]
\end{definition}

\begin{theorem}[Smooth-stratum duality]
If \(u:U\to S\) is smooth of pure relative dimension \(d\), then
\[
 \omega_{U,A}\xrightarrow{\sim}A_U(d)[2d].
\]
For every lisse-perfect \(K\) on \(U\),
\[
 D_{U,A}^{\mathrm{amb}}K\xrightarrow{\sim}
 K^{\vee_A}\otimes_A^L A_U(d)[2d],
 \qquad
 \delta_{U,A,K}:K\xrightarrow{\sim}D_{U,A}^{\mathrm{amb}}
 D_{U,A}^{\mathrm{amb}}K.
\]
These isomorphisms are natural in \(K\) and commute with restriction and
finite coefficient change.
\end{theorem}

\begin{theorem}[Existence of duality stratifications]
Because \(\bar k\) is perfect, every
\(K\in D_{\mathrm{ctf}}(X,A)\) admits a duality stratification.  Any two
such stratifications admit a common finite refinement by smooth locally
closed \(\bar k\)-subschemes.
\end{theorem}

\begin{theorem}[Open--closed localization and duality exchange]
Let \(j:U\hookrightarrow X\) be open and \(i:Z\hookrightarrow X\) its
closed complement.  The functor \(Ri^!\) is the right adjoint of the exact
functor \(i_*\).  For \(K\in D^b_c(X,A)\), there are distinguished
triangles
\[
 j_!j^*K\to K\to i_*i^*K\to j_!j^*K[1],
\qquad
 i_*Ri^!K\to K\to Rj_*j^*K\to i_*Ri^!K[1].
\]
On finite-Tor constructible objects there are natural isomorphisms
\[
 D_{X,A}j_!\xrightarrow{\sim}Rj_*D_{U,A},
 \qquad D_{X,A}Rj_*\xrightarrow{\sim}j_!D_{U,A},
\]
\[
 D_{X,A}i_*\xrightarrow{\sim}i_*D_{Z,A},
 \qquad i^*D_{X,A}\xrightarrow{\sim}D_{Z,A}Ri^!.
\]
They identify the dual of either localization triangle with the other one
and carry the three biduality maps to a morphism of distinguished triangles.
\end{theorem}

\begin{proposition}[One-stratum biduality]
If \(K\in D_{\mathrm{ctf}}(X,A)\) is supported on one smooth stratum and is
lisse perfect there, then \(D_{X,A}^{\mathrm{amb}}K\) is finite-Tor
constructible and \(\delta_{X,A,K}\) is an isomorphism.
\end{proposition}

\begin{proposition}[Localization induction for biduality]
In an open--closed localization triangle, suppose the first duals of the
open and closed restrictions are finite-Tor constructible and their
biduality morphisms are isomorphisms.  Then the first dual of the middle
object is finite-Tor constructible and its biduality morphism is an
isomorphism.
\end{proposition}

\begin{theorem}[Finite-stage stratified Verdier biduality]
For every \(n\ge1\) and every
\(K\in D_{\mathrm{ctf}}(X,A_n)\), one has
\[
 D_{X,A_n}K\in D_{\mathrm{ctf}}(X,A_n),
 \qquad
 \delta_{X,A_n,K}:K\xrightarrow{\sim}D_{X,A_n}D_{X,A_n}K.
\]
Thus
\[
 D_{X,A_n}:D_{\mathrm{ctf}}(X,A_n)^{\mathrm{op}}
       \xrightarrow{\sim}D_{\mathrm{ctf}}(X,A_n)
\]
is an exact anti-equivalence whose quasi-inverse is itself.
\end{theorem}

\begin{theorem}[Finite-stage duality under reduction]
For \(m\leq n\), the adjunction comparison and finite-coefficient
tensor--Hom base change give natural isomorphisms
\[
 \omega_{X,A_n}\otimes_{A_n}^LA_m\xrightarrow{\sim}\omega_{X,A_m},
 \qquad
 D_{X,A_n}K\otimes_{A_n}^LA_m
 \xrightarrow{\sim}D_{X,A_m}(K\otimes_{A_n}^LA_m).
\]
For \(l\leq m\leq n\), the direct comparison from \(A_n\) to \(A_l\)
equals the composite through \(A_m\), after the associativity isomorphism
for derived tensor product.  The comparisons commute with evaluation and
biduality.  Consequently \((\omega_{X,A_n})_n\) is a normalized tower and,
for every normalized finite-Tor tower \(K_\bullet\), the objects
\((D_{X,A_n}K_n)_n\) have canonical normalized transition maps.
\end{theorem}

\begin{definition}[Normalized finite-Tor adic duality category]
Let
\[
 D_{\mathrm{ctf}}(X,\mathcal O_E)_{\mathrm{adic}}
 \subseteq D^b_c(X,\mathcal O_E)_{\mathrm{adic}}
\]
be the full subcategory of towers \(K_\bullet\) for which there is one
interval \([r,s]\) such that every \(K_n\) has Tor amplitude in
\([r,s]\).
Define \(D_{\mathrm{ctf}}(S,\mathcal O_E)_{\mathrm{adic}}\) in the same
way.  The finite reduction isomorphisms define
\[
 \omega_{X,\mathcal O_E}=(\omega_{X,A_n})_n,
 \qquad
 D_{X,\mathcal O_E}(K_\bullet)
       =(D_{X,A_n}K_n)_n.
\]
Define
\[
 \omega_{X,E}
  =\operatorname{Real}(\omega_{X,\mathcal O_E})[1/\varpi_E],
 \qquad
 D_{X,E}K=R\mathcal Hom_E(K,\omega_{X,E}).
\]
\end{definition}

\begin{theorem}[Coefficient transport of Verdier duality]
The systems of dualizing objects, duals, evaluation maps, and biduality
maps are normalized.  Levelwise exceptional functors define an adjunction
of exact functors
\[
 Ra_{X,!}:D_{\mathrm{ctf}}(X,\mathcal O_E)_{\mathrm{adic}}
 \rightleftarrows D_{\mathrm{ctf}}(S,\mathcal O_E)_{\mathrm{adic}}:Ra_X^!,
\]
and \(D_{X,\mathcal O_E}\) is an exact anti-equivalence of the source
adic category with itself.  Adic realization followed by inversion of
\(\varpi_E\) gives
\[
 Ra_{X,!}:\Dbc(X,E)\rightleftarrows\Dbc(S,E):Ra_X^!,
 \qquad
 D_{X,E}:\Dbc(X,E)^{\mathrm{op}}\xrightarrow{\sim}\Dbc(X,E).
\]
For every finite extension
\(E\subset M\),
\[
 D_{X,E}K\otimes_E M\xrightarrow{\sim}D_{X,M}(K\otimes_E M),
\]
and these maps satisfy identity and transitivity.  Hence
\[
 D_X:\Dbc(X,E)^{\mathrm{op}}\xrightarrow{\sim}\Dbc(X,E),
 \qquad K\xrightarrow{\sim}D_XD_XK,
\]
with \(D_X=D_{X,E}\), and the same construction
descends to the Karoubi-completed finite-stage \(\Qlb\)-category.
\end{theorem}

\begin{theorem}[Compact-support Frobenius comparison]
The cartesian Frobenius square for \(a_X:X\to S\), together with
compactification independence and proper base change, defines a natural
isomorphism of exact functors
\[
 \gamma_{a,!}:\Phi_S^*Ra_{X,!}\xrightarrow{\sim}Ra_{X,!}\Phi_X^*.
\]
It is compatible with Frobenius iteration, the identity and composition
isomorphisms for compact support, finite reduction, adic realization,
rationalization, and finite scalar extension.
\end{theorem}

\begin{theorem}[Exceptional Frobenius exchange]
The right-adjoint mate of \(\gamma_{a,!}\) is an isomorphism
\[
 \gamma_{a,!}^{\,!}:\Phi_X^*Ra_X^!\xrightarrow{\sim}
        Ra_X^!\Phi_{S}^*.
\]
For the strict coefficient-unit arrow
\(\varphi_E:\Phi_S^*E\xrightarrow{\sim}E\), define
\[
 \beta_X:\Phi_X^*\omega_{X,E}
 \xrightarrow{\sim}Ra_X^!\Phi_S^*E
 \xrightarrow{Ra_X^!(\varphi_E)}Ra_X^!E=\omega_{X,E}.
\]
This arrow is compatible with iteration, the adjunction unit and counit,
finite reduction, rationalization, and finite scalar extension.
\end{theorem}

\begin{definition}[Strict Verdier dual]
For \((K,\varphi_K)\in\WeilD(X_0,E)\), define the strict Weil arrow on
\(D_XK=\RHom_E(K,\omega_{X,E})\) by
\[
 \varphi_{D_XK}:\Phi_X^*D_XK
 \xrightarrow{c^{\mathcal Hom}}
 \RHom_E(\Phi_X^*K,\Phi_X^*\omega_{X,E})
 \xrightarrow{\RHom_E(\varphi_K^{-1},\beta_X)}D_XK.
\]
\end{definition}

\begin{theorem}[Strict Verdier duality and equivariant biduality]
Strict Verdier duality is an exact anti-equivalence
\[
 D_X:\WeilD(X_0,E)^{\mathrm{op}}\xrightarrow{\sim}\WeilD(X_0,E).
\]
For every \((K,\varphi_K)\), the biduality isomorphism satisfies
\[
 \delta_K\circ\varphi_K
 =\varphi_{D_XD_XK}\circ\Phi_X^*(\delta_K).
\]
It commutes with finite scalar extension and extends to the finite-stage
\(\Qlb\)-category by
\[
                       D_X(E,K,e)=(E,D_XK,D_X(e)).
\]
\end{theorem}

\begin{theorem}[Finite-stage descent for Tate objects and Verdier duality]
Every finite collection of Tate objects, twists, dualizing complexes,
Verdier duals, strict Frobenius arrows, biduality maps, and coefficient
comparisons in \(\WeilD(X_0,\Qlb)_{\mathrm{2-colim}}\) is represented over
one common finite extension \(E/\Ql\).  Extension to a further finite coefficient-field
commutes with \(E(m)\), \(D_X\), biduality, and all displayed strict Weil
arrows.
\end{theorem}

\section{Finite Frobenius modules and spectra}

\begin{definition}[Finite Frobenius module]
A finite \(E\)-Frobenius module over the base cardinality \(Q>1\) is a
triple \((V,F;Q)\), where \(Q\) is an integer, \(V\) is a finite-dimensional
\(E\)-vector space, and \(F\in\operatorname{GL}_E(V)\).  The cardinality
\(Q\) is part of the datum, although it is omitted from notation in formulas
which involve only linear algebra.  At a closed point \(x\), it is always
\(Q=q_x\).
Its characteristic factor and characteristic polynomial are
\[
 P_{V,F}(T)=\det(1-TF),
 \qquad
 \chi_{V,F}(T)=\det(T-F).
\]
Its spectrum \(\operatorname{Sp}(V,F)\) is the multiset of roots of
\(\chi_{V,F}\) in \(\Qlb\), counted with algebraic multiplicity after a
finite splitting extension of \(E\).
For \(m\in\mathbf Z\), define the Tate twist of the full Frobenius datum by
\[
                   (V,F;Q)(m)=(V,Q^{-m}F;Q).
\]
Write \(V(m)\) for this full Frobenius module when \(F\) and \(Q\) are
fixed by the context.
\end{definition}

\begin{proposition}[Exact sequences and spectra]
For a Frobenius-equivariant short exact sequence of modules over the same
base cardinality \(Q\)
\[
             0\longrightarrow V'\longrightarrow V\longrightarrow V''
             \longrightarrow0
\]
one has
\[
 P_V(T)=P_{V'}(T)P_{V''}(T),\qquad
 \operatorname{Sp}(V)=\operatorname{Sp}(V')\sqcup
                       \operatorname{Sp}(V'').
\]
The spectrum of a Frobenius-stable subquotient is a submultiset of the
spectrum of \(V\).  Consequently every eigenvalue condition stable under
submultisets is stable under Frobenius-stable subobjects and quotients.
\end{proposition}

\begin{proposition}[Spectral operations]
For finite Frobenius modules \((V,F_V;Q)\), \((W,F_W;Q)\) over the same
base cardinality and \(m\in\mathbf Z\), Frobenius on
\(V\otimes W\), \(V^\vee\), and \(V(m)\) is respectively
\[
 F_V\otimes F_W,\qquad (F_V^{-1})^\vee,\qquad Q^{-m}F_V
\]
respectively.  Their spectra are
\[
 \{\alpha\beta\},\qquad\{\alpha^{-1}\},\qquad\{Q^{-m}\alpha\},
\]
with product multiplicities in the tensor case.
\end{proposition}

\section{Scalar weights}

\begin{definition}[All-conjugates Weil number]
Let \(Q>1\) be an integer and \(w\in\mathbf R\).  A nonzero
\(\alpha\in\Qlb\) is a \(Q\)-Weil number of weight \(w\) if it is
algebraic over \(\mathbf Q\) and
\[
                  |\sigma(\alpha)|=Q^{w/2}
\]
for every embedding
\(\sigma:\mathbf Q(\alpha)\hookrightarrow\mathbf C\).  Algebraicity and
the quantification over every complex conjugate are part of the definition.
\end{definition}

\begin{definition}[Fixed-\(\iota\) scalar weight and bounds]
For nonzero \(\alpha\in\Qlb\), define
\[
 \operatorname{wt}_{\iota,Q}(\alpha)=w
 \quad\Longleftrightarrow\quad
 |\iota(\alpha)|=Q^{w/2}.
\]
Define weight at most \(w\) and weight at least \(w\) by replacing equality
with \(\le\) and \(\ge\), respectively.  Because \(Q>1\), a fixed-\(\iota\)
weight, when it exists, is unique.
\end{definition}

\begin{definition}[Q-indexed weights of a finite Frobenius module]
Let \((V,F;Q)\) be a finite \(E\)-Frobenius module and let
\(w\in\mathbf R\).  Define
\[
\begin{aligned}
 \operatorname{Pure}_{\iota,Q}(V,F;w)
 &\Longleftrightarrow
   \operatorname{wt}_{\iota,Q}(\alpha)=w
   \quad\text{for every }\alpha\in\operatorname{Sp}(V,F),\\
 \operatorname{WeilPure}_{Q}(V,F;w)
 &\Longleftrightarrow
   \alpha\text{ is a \(Q\)-Weil number of weight \(w\)}
   \quad\text{for every }\alpha\in\operatorname{Sp}(V,F),\\
 \operatorname{WtLe}_{\iota,Q}(V,F;w)
 &\Longleftrightarrow
   |\iota(\alpha)|\le Q^{w/2}
   \quad\text{for every }\alpha\in\operatorname{Sp}(V,F),\\
 \operatorname{WtGe}_{\iota,Q}(V,F;w)
 &\Longleftrightarrow
   |\iota(\alpha)|\ge Q^{w/2}
   \quad\text{for every }\alpha\in\operatorname{Sp}(V,F).
\end{aligned}
\]
The predicate \(\operatorname{IntMixLe}_{\iota,Q}(V,F;w)\) means that,
after a finite splitting extension of \(E\), the module has a finite
Frobenius-stable filtration whose nonzero quotients satisfy
\(\operatorname{Pure}_{\iota,Q}(-;u)\) for integers \(u\le w\).  Define
\(\operatorname{IntMixGe}_{\iota,Q}(V,F;w)\) by integers \(u\ge w\).
The displayed parameter \(Q\) is retained under scalar extension.
\end{definition}

\begin{proposition}[Power change of the base cardinality]
Let \(r\ge1\).  The assignment
\[
                 (V,F;Q)\longmapsto(V,F^r;Q^r)
\]
preserves and reflects
\(\operatorname{Pure}_{\iota,Q}(-;w)\),
\(\operatorname{WeilPure}_{Q}(-;w)\),
\(\operatorname{WtLe}_{\iota,Q}(-;w)\),
\(\operatorname{WtGe}_{\iota,Q}(-;w)\),
\(\operatorname{IntMixLe}_{\iota,Q}(-;w)\), and
\(\operatorname{IntMixGe}_{\iota,Q}(-;w)\), with \(Q\) replaced by
\(Q^r\) on the target side.  The assertion includes inseparable-free
finite base-field extension, where local Frobenius and residue cardinality
are raised to the same positive power.
\end{proposition}

\begin{proposition}[Scalar weight algebra]
If \(\alpha\) and \(\beta\) have fixed-\(\iota\) weights \(u\) and \(v\)
relative to the same \(Q\), then
\[
 \operatorname{wt}_{\iota,Q}(\alpha\beta)=u+v,
 \qquad \operatorname{wt}_{\iota,Q}(\alpha^{-1})=-u,
 \qquad \operatorname{wt}_{\iota,Q}(\alpha^n)=nu.
\]
If \(Q=q_x=q^{d(x)}\), multiplication by \(q_x^{-m}\) changes the weight
from \(u\) to \(u-2m\).  The same statements hold for all-conjugates Weil
numbers.
\end{proposition}

\begin{theorem}[All-conjugates adapter]
Let \(Q>1\) be an integer, let \(w\in\mathbf R\), and let
\(\alpha\in\Qlb^\times\) be algebraic over \(\mathbf Q\).  The following
conditions are equivalent.
\begin{enumerate}
\item \(\alpha\) is a \(Q\)-Weil number of weight \(w\).
\item For every embedding
\(\sigma:\mathbf Q(\alpha)\hookrightarrow\mathbf C\), one has
\(|\sigma(\alpha)|=Q^{w/2}\).
\item For every \(\tau\in\operatorname{Aut}(\Qlb/\mathbf Q)\), the scalar
\(\tau(\alpha)\) has fixed-\(\iota\) weight \(w\).
\end{enumerate}
In particular, all-conjugates purity implies fixed-\(\iota\) purity for
every choice of \(\iota\).  Purity for one fixed embedding alone does not
imply the all-conjugates condition.
\end{theorem}

\section{Punctual purity and mixed sheaves}

\begin{definition}[Closed-point spectrum]
For a constructible Weil sheaf \(\mathcal F\) and \(x\in|X_0|\), define
\[
 \operatorname{Fib}_x(\mathcal F)
   =(\mathcal F_{\bar x},F_x;q_x),
 \qquad
        \operatorname{Sp}_x(\mathcal F)
        =\operatorname{Sp}(\mathcal F_{\bar x},F_x).
\]
It is independent of \(\bar x\) and is unchanged by enlarging the finite
coefficient field except for viewing its elements in the larger field.
\end{definition}

\begin{definition}[Punctual purity]
A constructible Weil sheaf \(\mathcal F\) is punctually
\(\iota\)-pure of weight \(w\in\mathbf R\) if, for every closed point \(x\)
one has
\(\operatorname{Pure}_{\iota,q_x}(\mathcal F_{\bar x},F_x;w)\).
It is punctually Weil-pure of weight \(w\) if
\(\operatorname{WeilPure}_{q_x}(\mathcal F_{\bar x},F_x;w)\) holds for
every closed point \(x\).  The first is a fixed-embedding condition; the
second is an all-conjugates condition.
\end{definition}

\begin{definition}[Mixed sheaf and sheafwise bounds]
A constructible \(E\)-Weil sheaf \(\mathcal F\) is mixed if there is a finite
extension \(M/E\) inside \(\Qlb\) and a finite increasing filtration by
\(M\)-Weil subsheaves
\[
 0=W_{-1}\mathcal F_M\subset W_0\mathcal F_M\subset\cdots
   \subset W_N\mathcal F_M=\mathcal F_M,
 \qquad \mathcal F_M=\mathcal F\otimes_E M,
\]
such that every nonzero quotient
\(W_a\mathcal F_M/W_{a-1}\mathcal F_M\) is punctually
\(\iota\)-pure of some real weight \(w_a\).  It has weights at most \(w\)
if all \(w_a\le w\), and weights at least \(w\) if all \(w_a\ge w\).
The numerical index \(a\) labels the filtration steps and is not itself the
weight.
The sheaf is \emph{integrally mixed} if every \(w_a\) can be chosen in
\(\mathbf Z\).  Write \(\operatorname{IntWtLe}(\mathcal F,w)\),
respectively \(\operatorname{IntWtGe}(\mathcal F,w)\), when such an
integral filtration can be chosen with every \(w_a\le w\), respectively
every \(w_a\ge w\).
\end{definition}

\begin{theorem}[Abelian stability of mixed sheaves]
Mixed constructible Weil sheaves form an abelian subcategory of
\(\mathsf{Cons}_{\mathrm W}(X_0,E)\).  A fixed upper or lower sheafwise
bound is stable under subobjects, quotients, extensions, finite direct sums,
and direct summands.  These assertions are preserved and reflected by
finite coefficient extension.

\end{theorem}

\begin{proposition}[Weight separation for sheaves]
Let \(\mathcal F\) be punctually \(\iota\)-pure of weight \(u\), and let
\(\mathcal G\) be punctually \(\iota\)-pure of weight \(v\), with
\(u\ne v\).  Every
Weil-sheaf morphism \(\mathcal F\to\mathcal G\) is zero.  More generally,
if \(\mathcal F\) has weights at most \(u\) and \(\mathcal G\) has weights
at least \(v>u\), then
\[
             \operatorname{Hom}_{\mathsf{Cons}_{\mathrm W}}
             (\mathcal F,\mathcal G)=0.
\]
\end{proposition}

\begin{theorem}[Scalar and base-field invariance]
Let \(M/E\) be finite.  A constructible \(E\)-Weil sheaf is punctually
\(\iota\)-pure,
mixed of weights at most \(w\), or mixed of weights at least \(w\) if and
only if its scalar extension to \(M\) has the corresponding property.

Let \(k_r=\mathbf F_{q^r}\), and let \(\mathcal F_{k_r}\) be the pullback of
\(\mathcal F\) after base-field extension.  At a closed point above \(x\),
the new geometric Frobenius is a positive power of \(F_x\), while the new
residue-field cardinality is the corresponding power of \(q_x\).  Hence
punctual purity and both mixed bounds are preserved and reflected by finite
extension of the base field.
\end{theorem}

\section{BBD weights on complexes}

\begin{definition}[Upper complex weight]
For \(K\in\WeilD(X_0,E)\) and \(w\in\mathbf R\), define
\[
\WtLe(K,w)
 \quad\Longleftrightarrow\quad
 \mathcal H^i(K)\text{ is mixed of weights at most }w+i
 \text{ for every }i\in\mathbf Z.
\]
Define \(\operatorname{IntWtLe}(K,w)\) by requiring every
\(\mathcal H^i(K)\) to be integrally mixed of weights at most \(w+i\).
Only finitely many conditions are nontrivial because \(K\) is bounded.
\end{definition}

\begin{definition}[Lower complex weight and purity]
Using the strict Verdier anti-equivalence constructed above, define the BBD
lower predicates by
\[
             \WtGe(K,w)\quad\Longleftrightarrow\quad
             \WtLe(D_XK,-w).
\]
Define
\[
 \operatorname{IntWtGe}(K,w)\quad\Longleftrightarrow\quad
 \operatorname{IntWtLe}(D_XK,-w).
\]
The complex \(K\) is pure of weight \(w\) if both \(\WtLe(K,w)\) and
\(\WtGe(K,w)\) hold.  It has weights in \([a,b]\) if
\(\WtGe(K,a)\) and \(\WtLe(K,b)\) both hold.
\end{definition}

\begin{definition}[Weights on finite-stage retracts]
For a retract \((E,K,e)\) in the finite-stage \(\Qlb\)-category, apply the
upper definitions to the images of the cohomology idempotents and apply the
lower definitions to \((E,D_XK,D_Xe)\).
\end{definition}

\begin{proposition}[Independence of a finite-stage representative]
The four upper and lower predicates on a retract are unchanged by extension
to a common finite coefficient field and by an isomorphism of idempotent
representatives.  They therefore define predicates on
\(\WeilD(X_0,\Qlb)_{\mathrm{2-colim}}\).
\end{proposition}

\begin{proposition}[The point case]
On \(\Spec\Fq\), Verdier duality is the ordinary linear dual with
inverse-transpose Frobenius.  Consequently
\[
 \WtGe(K,w)
 \Longleftrightarrow
 \mathcal H^i(K)\text{ has weights at least }w+i
 \text{ for every }i.
\]
On a general singular scheme the lower predicate remains the Verdier-dual
condition and is not replaced by this pointwise cohomology formula.
\end{proposition}

\begin{proposition}[Shifts, twists, and duality]
For \(m,n\in\mathbf Z\),
\[
\begin{aligned}
 \WtLe(K,w)&\Longleftrightarrow\WtLe(K[n],w+n),&
 \WtGe(K,w)&\Longleftrightarrow\WtGe(K[n],w+n),\\
 \WtLe(K,w)&\Longleftrightarrow\WtLe(K(m),w-2m),&
 \WtGe(K,w)&\Longleftrightarrow\WtGe(K(m),w-2m).
\end{aligned}
\]
Moreover
\[
 \WtLe(K,w)\Longleftrightarrow\WtGe(D_XK,-w),
 \qquad
 \WtGe(K,w)\Longleftrightarrow\WtLe(D_XK,-w).
\]
Thus \(D_X\), \([n]\), and \((m)\) send a pure weight \(w\) respectively
to the pure weights \(-w\), \(w+n\), and \(w-2m\).
Every displayed equivalence remains valid after replacing
\(\WtLe,\WtGe\) by \(\operatorname{IntWtLe},\operatorname{IntWtGe}\).
\end{proposition}

\begin{proposition}[Correctly shifted triangle rules]
Let
\[
                 A\longrightarrow B\longrightarrow C\longrightarrow A[1]
\]
be a distinguished triangle.  The upper rules are
\[
\begin{array}{lll}
 \WtLe(A,a),\ \WtLe(C,c)&\Longrightarrow&
       \WtLe(B,\max\{a,c\}),\\
 \WtLe(A,a),\ \WtLe(B,b)&\Longrightarrow&
       \WtLe(C,\max\{a+1,b\}),\\
 \WtLe(B,b),\ \WtLe(C,c)&\Longrightarrow&
       \WtLe(A,\max\{b,c-1\}).
\end{array}
\]
The lower rules are
\[
\begin{array}{lll}
 \WtGe(A,a),\ \WtGe(C,c)&\Longrightarrow&
       \WtGe(B,\min\{a,c\}),\\
 \WtGe(A,a),\ \WtGe(B,b)&\Longrightarrow&
       \WtGe(C,\min\{a+1,b\}),\\
 \WtGe(B,b),\ \WtGe(C,c)&\Longrightarrow&
       \WtGe(A,\min\{b,c-1\}).
\end{array}
\]
In particular, a fixed upper bound and a fixed lower bound are each stable
under extensions, meaning the first rule in the corresponding display.
The same six implications hold for the integral predicates
\(\operatorname{IntWtLe}\) and \(\operatorname{IntWtGe}\).
\end{proposition}

\begin{proposition}[Finite sums, retracts, and finite stages]
The upper bound of a finite direct sum is bounded by the maximum of the
upper bounds of its summands, and its lower bound is bounded by the minimum
of their lower bounds.  Every bound on an object is inherited by a retract.
All weight predicates are preserved and reflected by finite scalar
extension and therefore define predicates on
\(\WeilD(X_0,\Qlb)_{\mathrm{2-colim}}\).
\end{proposition}

\begin{proposition}[Frobenius spectral sequences]
Let \((E_r^{a,b},F_r^{a,b};Q)\) be a bounded spectral sequence of finite
Frobenius modules over one fixed base cardinality \(Q>1\), whose
differentials commute with Frobenius, and suppose the abutment
\((H^n,F^n;Q)\) has a finite Frobenius-stable filtration with graded pieces
\(E_\infty^{a,n-a}\).  If
\[
 \operatorname{WtLe}_{\iota,Q}
   (E_1^{a,b},F_1^{a,b};w+a+b)
 \quad\text{for every }a,b,
\]
then
\(\operatorname{WtLe}_{\iota,Q}(H^n,F^n;w+n)\).  If instead every
\(E_1^{a,b}\) satisfies
\(\operatorname{WtGe}_{\iota,Q}(-;w+a+b)\), then
\(\operatorname{WtGe}_{\iota,Q}(H^n,F^n;w+n)\).
\end{proposition}

\section{Nilpotent monodromy}

\begin{definition}[Nilpotent operator and monodromy filtration]
Let \(V\) be a finite-dimensional \(E\)-vector space and let
\(N:V\to V\) be nilpotent.  A monodromy filtration centered at zero is a
finite increasing filtration \((M_jV)_{j\in\mathbf Z}\) such that
\[
 N(M_jV)\subset M_{j-2}V,
 \qquad
 N^j:\gr_j^M V\xrightarrow{\sim}\gr_{-j}^M V
 \quad(j\ge0),
\]
where \(\gr_j^M V=M_jV/M_{j-1}V\).  A filtration centered at \(c\in\mathbf Z\)
is obtained by reindexing \(M[c]_j=M_{j-c}\).
\end{definition}

\begin{theorem}[Existence and uniqueness of the monodromy filtration]
Every nilpotent endomorphism of a finite-dimensional vector space has a
unique monodromy filtration centered at zero.  If a Jordan chain has length
\(d\), and its basis is ordered by
\[
 N e_0=0,\qquad N e_t=e_{t-1}\quad(1\le t<d),
\]
then \(e_t\) has filtration degree \(2t-(d-1)\).  The filtration obtained by
spanning all basis vectors of degree at most \(j\) is independent of the
chosen Jordan decomposition.

\end{theorem}

\begin{definition}[Primitive pieces]
For \(j\ge0\), define
\[
 P_j(V,N)=\ker\!\left(
 N^{j+1}:\gr_j^M V\longrightarrow\gr_{-j-2}^M V
\right).
\]
For \(j<0\), define \(P_j(V,N)=0\).  With this convention the
primitive-decomposition map in graded degree \(r\in\mathbf Z\) is
\[
       \bigoplus_{a\ge0}N^aP_{r+2a}(V,N)
       \longrightarrow\gr_r^M V.
\]
Only finitely many summands are nonzero.
\end{definition}

\begin{theorem}[Primitive decomposition]
For every nilpotent pair \((V,N)\), the primitive-decomposition map is an
isomorphism in every degree.  It is natural for isomorphisms of nilpotent
pairs and is independent of any Jordan or \(\mathfrak{sl}_2\)-splitting.
\end{theorem}

\begin{proposition}[Tensor and dual monodromy]
For nilpotent pairs \((V,N_V)\) and \((W,N_W)\), put
\[
 N_{V\otimes W}=N_V\otimes1+1\otimes N_W,
 \qquad
 (N_{V^\vee}\lambda)(v)=-\lambda(N_Vv).
\]
Their monodromy filtrations are characterized by
\[
 M_n(V\otimes W)=\sum_{a+b=n}M_aV\otimes M_bW,
 \qquad
 M_j(V^\vee)=\{\lambda:\lambda(M_{-j-1}V)=0\}.
\]
The corresponding graded maps identify
\(\gr_j^M(V^\vee)\simeq(\gr_{-j}^M V)^\vee\).
\end{proposition}

\begin{definition}[Weil--Deligne weight]
A finite Weil--Deligne module over a base cardinality \(Q>1\) is a
quadruple \((V,F,N;Q)\), where \((V,F;Q)\) is a finite Frobenius module,
\(N\) is nilpotent, and
\[
                         FNF^{-1}=Q^{-1}N.
\]
Its monodromy filtration is Frobenius-stable.  It is \(\iota\)-pure of
Weil--Deligne weight \(w\) if
\[
 \operatorname{Pure}_{\iota,Q}
       (\gr_j^M V,\gr_j^M F;w+j)
 \quad\text{for every }j.
\]
The relation with \(F\) makes \(N\) lower Frobenius weight by two, as
required by the filtration shift \(j\mapsto j-2\).
\end{definition}

\end{document}
